\section{Data description}

In our analysis, we used second and third waves (NLFS 2008/09 and NLFS 2017/18) of \textit{Nepal Labor Force Survey} administrated by the National Statistics Office (NSO)(formerly known as Central Bureau of Statistics(CBS)) in collaboration with International Labor Organization (ILO) Nepal. Using a multi-stage clustering sampling approach, these survey collects rich information about individual labor pertaining to demographic profiles, labor market indicators, employment and unemployment and migration. NLFS 2008/09 (referred to as NLFS II hereafter) included interviews with 16,000 households, while the NLFS 2017/18 (referred to as NLFS III) expanded to 18,000 households. We restrict our analysis to individuals of 15-65 years who are in full time wage employment defined as individuals who work for 40 hours or more during the week. We did not include agricultural job sector in the sample, in accordance with This resulted in working-age samples of 45,222 individuals in the second round and 47,899 individuals in the third. From this, we classified employed samples comprising 6718 individuals (35\% formal and 65\% informal employment) and 6827 individuals (23\% formal and 77\% informal employment) in the NLFS II and III respectively. 

We use the measure of job-based concept of informal employment within the working-age samples. For this, we consider the non-agricultural sector  We identify employees in the formal employment as those who work in government service, financial public corporation, non-financial public corporation, governmental and non-governmental organisations and private registered financial and non-financial companies, have a permanent employment and their employer pays social security contribution or they receive the benefits of paid leave or get compensation for unused leave. Based on the same, we identify informal employment by exclusion, where we consider those not falling into these category as informal employment. 

The dependent variable in our study is the logarithm of the hourly income. The survey collects information on the usual hours work per week but not the number of weeks worked during a month. We compute hourly income by dividing monthly income by total imputed hourly work per month. Therefore, we compute the monthly hours of work by multiplying the usual hours of work per week to four. The monthly income is net income paid to people. Thus it includes transfers such as pension, social insurance deductions, and taxes as well as extra income like bonuses, overtime. premiums etc.

We investigate the determinants of wage in our study through the lens of examining human capitaL variables, personal and household, job and employer characteristics. We focus on experience, experience squared, and education as human capital variables, as extensive literature highlights their significant association with wage levels (Becker, 1975; Borjas, 2005; Rosen, 1972; Mincer, 1974), including the number of years with an employer (Burdett and Cole, 2010). Personal and household characteristics such as gender, caste group, marital status, and the number of children are included, given their proven impact on wages (Bardasi and Taylor, 2008; Killewald and Gough, 2013; Becker, 1975). Job characteristics encompass occupation class, vocational training, migration for work, overtime, and total household chores, all identified as influential in previous research (Becker, 1975; Borjas, 2005). Employer characteristics are captured through the job sector by economic sectors, reflecting their significant role in wage determination (Becker, 1975; Borjas, 2005).

We identity overtime work as someone 

%The 17th International Conference of Labour Statisticians (ICLS) held by the International Labour Organization (ILO) in 2003 acknowledged the existence of informal employment within both formal and informal enterprises, noting that the traditional enterprise-based definition overlooked significant aspects of informality. Consequently, the 17th ICLS introduced a conceptual framework that aligned the enterprise-based concept of employment in the informal sector with a job-based concept of informal employment, thereby broadening the scope. Informal employment was defined as jobs not subject to standard labour legislation, taxation, social or legal protections to certain employment benefits regardless of whether employed by formal or informal sector enterprises.% 

